\TieudeT{PHONG TỎA VẬT LÍ}
\subsubsection*{PHẦN I. Thí sinh trả lời câu hỏi từ câu 1 đến câu 18. Mỗi câu hỏi thí sinh chỉ chọn một phương án}
\setcounter{ex}{0}
\begin{ex}%book%[1P1K5-3][Loai1] 
Chất điểm chuyển động dọc theo trục $Ox$ với tốc độ $3,5\ m/s$ thì tăng tốc chuyển động nhanh dần đều, trong $2\ s$ tốc độ tăng đến $4,5\ m/s$. Quãng đường trong thời gian $2\ s$ đó là
\choice
{$5\ m$}
{$6\ m$}
{$7\ m$}
{\True $8\ m$}
\haidong{5}
\loigiai{
\begin{itemize}
\item Ta có: $v_t=v_0+at\Leftrightarrow 4,5=3,5+2a\Rightarrow a=0,5\ m/s^2$.
\item Lại có: $s=\dfrac{v_t^2-v_0^2}{2a}=8\ m$.
\end{itemize}
}
\haidong{7} 
\end{ex}
\begin{ex}%book %[1P1-5.2B][Loai1] 
Một xe máy đang đi với tốc độ $36\ km/h$ bỗng người lái xe thấy có một cái hố trước mặt, cách xe $20\ m$. Người ấy phanh gấp và xe đến sát miệng hố thì dừng lại. Thời gian hãm phanh là
\choice
{$2\ s$}
{$0,2\ s$}
{$0,4\ s$}
{\True $4\ s$}
\haidong{5}
\loigiai{
\begin{itemize}
\item Đổi $36\  km/h = 10\  m/s$.
\item Gia tốc của chuyển động là: $a=\dfrac{v^2-v_0^2}{2d}=\dfrac{0^2-10^2}{2.20}=- 2,5 \  m/s^2$.
\item Nên thời gian hãm phanh là: $t=\dfrac{v-v_0}{a}=\dfrac{0-10}{-2,5}=4\ s$.
\end{itemize}	
}
\haidong{7} \end{ex}
\begin{ex}%book%[1P1-5.3K][Loai1]
Một đoàn tàu đang đi với tốc độ $10\ m/s$ thì hãm phanh, chuyển động chậm dần đều. Sau khi đi thêm được $64\ m$ thì tốc độ của nó chỉ còn $21,6\ km/h$. Quãng đường tàu đi thêm được kể từ lúc hãm phanh đến lúc dừng lại là
\choice
{$s = 120\ m$}
{$s = 110\ m$}
{\True $s = 100\ m$}
{$s = 200\ m$}	
\haidong{4}
\loigiai{
\begin{itemize}
\item Đổi $21,6\ km/h = 6\ m/s$.
\item Ta có: 
\begin{align*}
&v =v_0+at \Rightarrow v^2-v_0^2=\left(v_0+at\right)^2-v_0^2=2a\left(\dfrac{at^2}{2}+v_0t\right)\\ \Leftrightarrow &v^2-v_0^2=2as \Rightarrow 6^2-10^2=2a.64 \Rightarrow a=-0,5\ m/s^2.
\end{align*}
\item Quãng đường xe đi được từ lúc hãm phanh đến lúc dừng lại là: $s=\dfrac{v^2-v_0^2}{2a}=\dfrac{0-10^2}{2.\left(-0,5\right)}=100\ m$. 		
\end{itemize}	
}
\haidong{7} \end{ex}
\begin{ex}%book%[1P1K5-2][Loai1]
Một đoàn tàu đang chạy với tốc độ $54\ km/h$ thì hãm phanh và chuyển động chậm dần đều, sau đó đi thêm $125 m$ nữa thì dừng hẳn. Sau $5$ giây từ khi hãm phanh thì tàu chạy với tốc độ bằng
\choice
{$7,5\ m/s$}
{\True $10,5\ m/s$}
{$15\ m/s$}
{$5\ m/s$}	
\haidong{4}
\loigiai{
\begin{itemize}
\item $54\ km/h = 15\ m/s$.
\item Gia tốc đoàn tàu là: $a=\dfrac{0-15^2}{2.125}=0,9\ m/s^2$. 
\item Vận tốc đoàn tàu sau 5 s hãm phanh là: $v_t = 15 -0,9.5 = 10,5\ m/s$.
\end{itemize}	
}
\haidong{7} \end{ex}
\begin{ex}%book%8.5
Một chiếc ô tô đang chạy với vận tốc $23\ m/s$ thì đạp phanh. Sau $10\ s$, vận tốc của ô tô chỉ còn $11\ m/s$. Gia tốc của ô tô có giá trị đại số là
\choice
{$a=-1,2\ m/s^2$}
{$a=12\ m/s^2$}
{$a=1,2\ m/s^2$}
{$a=-12\ m/s^2$}
\haidong{7}
\loigiai{
\begin{itemize}
\item $a=-1,2\ m/s^2$
\item Gia tốc a có giá trị âm, các vận tốc có giá trị dương vì chuyển động là chậm dần.
\end{itemize}	
}
\haidong{7} \end{ex}
\begin{ex}%book%[1P1-5.3K][Loai1] 
Xe máy chuyển động thẳng nhanh dần đều từ trạng thái nghỉ, trong $3\ s$ đầu đi được quãng đường $2,5\ m$. Quãng đường xe máy đã đi trong giây thứ $3$ là
\choice
{\True $\dfrac{25}{18}\ m$}
{$\dfrac{25}{7}\ m$}
{$3\ m$}
{$\dfrac{9}{7}\ m$}
\haidong{7}
\loigiai{
\begin{itemize}
\item Ta có: $s=v_0t+\dfrac{a}{2}{t}^2$. 
\item Vì xe chuyển động nhanh dần từ trạng thái nghỉ nên $v_0=0$.
\item Ta có: $s=\dfrac{1}{2}a{t}^2$.
\item Quãng đường xe máy đi được trong 2 giây đầu tiên: $s_{t=2}=\dfrac{1}{2}a{t}^2=\dfrac{1}{2}.\dfrac{5}{9}{.2}^2=\dfrac{10}{9}\ m$.
\item Quãng đường xe máy đi được trong 3 giây đầu tiên: $s_{t=3}=\dfrac{1}{2}a{t}^2=\dfrac{1}{2}.\dfrac{5}{9}{.3}^2=\dfrac{5}{2}\ m$.
\item Vậy quãng đường đi được trong giây thứ 3 là: $\Delta s=s_{(t=3)}-s_{(t=2)}=\dfrac{25}{18}\ m$.
\end{itemize}
}
\haidong{7} \end{ex}
\begin{ex}%book%[1P1G5-3][Loai1] 
Xe máy chuyển động thẳng nhanh dần đều, trong $5\ s$ đầu đi được quãng đường $8,75\ m$. Biết vận tốc xe máy lúc $t = 3\ s$ là $v_t = 2\ m/s$. Quãng đường xe máy đi trong $10\ s$ kể từ cuối giây thứ $5$ là
\choice
{$45\ m$}
{$50\ m$}
{\True $55\ m$}
{$60\ m$}
\haidong{9}
\loigiai{
\begin{itemize}
\item Ta có: $\begin{cases}
s=v_0t+\dfrac{a}{2}t^2\\ 
v_t=v_0+at
\end{cases}\Rightarrow \begin{cases}
8,75=5v_0+12,5a\\ 
2=v_0+3a 
\end{cases}\Rightarrow \left\{{\begin{matrix}
{v_0=0,5\ m/s}\\
{a=0,5\ m/s^2}\\
\end{matrix}}\right.$
\item Quãng đường xe máy đi được trong 5 s: $s_{t=5}=v_0t+\dfrac{a}{2}{t}^2=8,75\ m$.
\item Quãng đường xe máy đi được trong 15 s: $s_{t=15}=v_0t+\dfrac{a}{2}{t}^2=63,75\ m$.
\item Quãng đường đi được trong 10s kể từ cuối giây thứ 5: $\Delta s=s_{t=15}-s_{t=5}=55\ m$.
\end{itemize}
}	
\haidong{7} \end{ex}
\begin{ex}%Tailieuchuan
Xe ô tô đang chuyển động thẳng với vận tốc $20 \ m/s$ thì bị hãm phanh chuyển động chậm dần đều. Quãng đường xe đi được từ lúc hãm phanh đến khi xe dừng hẳn là $100\ m$. Gia tốc của xe có độ lớn là
\choice
{$1 \ m/s^2$}
{$1.5 \ m/s^2$}
{\True $2 \ m/s^2$}
{$2,55 \ m/s^2$}
\loigiai{}
\end{ex}
\begin{ex}%book%[1P1-6.2K][Loai1]
\immini[thm]{Một xe máy chuyển động thẳng biến đổi đều có đồ thị vận tốc - thời gian như hình vẽ. Quãng đường đi được trong $20\ s$ đầu tiên là
\choice[2]
{\True $120\ m$}
{$100\ m$}
{$80\ m$}
{$90\ m$}	
}{\begin{hinhanh}{6}\begin{tikzpicture}[scale=1,thick,>=stealth',x=1cm,y=0.06cm,draw=Mapcolor]
\tkzDefPoints{0/0/o, 0/50/x, 5/0/t, 0/20/m, 4/40/n, 4/0/n1, 0/40/n2}
\tkzDrawSegments[dashed](n2,n n1,n)
\tkzDrawLines[very thick,color=Mapcolor,add = 0 and 0.2](o,n)
\tkzDrawSegments[->](o,x o,t)
\tkzLabelPoint[below](o){$O$}
\tkzLabelPoint[left](n2){$12$}
\tkzLabelPoint[below](n1){$20$}
\tkzLabelPoint[right](x){$v\ m/s$}
\tkzLabelPoint[above](t){$t\ (s)$}
\end{tikzpicture}\end{hinhanh}}
\haidong{9}
\loigiai{
Từ đồ thị ta có: $\begin{cases}
v_0=0\\
t=20\ s\\
\end{cases}\Rightarrow a=\dfrac{v-v_0}{t_{OA}}=\dfrac{12-0}{20}=0,6\ m/s^2$.\\
Quãng đường đi được là: $s=v_0t+\dfrac{1}{2}at^2=\dfrac{1}{2}.0,6.20^2=120\ m$.
}
\haidong{10} \end{ex}
\begin{ex}%book
\immini[thm]{Một vật chuyển động thẳng có đồ thị vận tốc theo thời gian như hình vẽ bên cạnh. Quãng đường vật đó đi được kể từ thời điểm $0\ s$ đến giây thứ $4$ là
\choice[2]
{$7\ m$}
{$8\ m$}
{\True $11\ m$}
{$4\ m$}}{\begin{hinhanh}{5}\begin{tikzpicture}[very thick,draw=Mapcolor]
\tkzDefPoints{0/0/o, 0/3/v, 4/0/t, 1.4/2/a, 2.1/2/b, 2.8/1/c}			
\draw  [help  lines, dashed, xstep=0.7cm,ystep=0.5cm]  (0,0)  grid  (3.8,2.5);
\tkzDrawSegments[Mapcolor, very thick](o,a a,b b,c)
\tkzDrawSegments[->](o,v o,t)
\node [left] at (0,2) {4};

\tkzLabelPoint[below](o){$O$}
\tkzLabelPoint[above](v){$v~(m/s)$}
\tkzLabelPoint[below](t){$t~(s)$}
\node [below] at (1.4,0) {2};
\node [below] at (2.8,0) {4};
\end{tikzpicture}\end{hinhanh}}
\haidong{6}
\loigiai{
}
\haidong{10} \end{ex}
\begin{ex}%Tailieuchuan
\immini[thm]{Cho đồ thị có dạng như hình dưới. Đây là chuyển động
\choice
{thẳng đều}
{nhanh dần đều}
{\True chậm dần đều}
{tròn đều}}{\begin{hinhanh}{6}\begin{tikzpicture}[draw=Mapcolor, very thick,line cap=round, line join=round, scale=1, font=\footnotesize, >=stealth,draw=Mapcolor]
\path
(0,0)coordinate(O)
(4,0)coordinate(x)
(0,4)coordinate(y)
;

\draw [->](O)--(x);
\draw [->](O)--(y);

\node at (0,4.5) {$v$};
\node at (4.5,0) {$t$};
\node[below left] at (0,0) {$O$};

\foreach  \x  in  {1,2,3}
{\pgfmathtruncatemacro{\label}{\x};
\draw[fill=white](\x,0)circle(0.5pt);}

\foreach  \y  in  {1,2,3}
{\pgfmathtruncatemacro{\label}{10*\y};
\draw[fill=white](0,\y)circle(0.5pt);}


\draw[very thick,Mapcolor] (0,3) -- (3,1);
\end{tikzpicture}\end{hinhanh}
}
\loigiai{

}
\end{ex}
\begin{ex}%Tailieuchuan
Một ô tô chuyển động thẳng biến đổi đều từ trạng thái nghỉ, đạt vận tốc $20 \ m/s$ sau $5\ s$. Quãng đường mà ô tô đã đi được là
\choice
{$100\ m$}
{\True $50\ m$}
{$25\ m$}
{$200\ m$}
\loigiai{

}
\end{ex}

\begin{ex}%Tailieuchuan
Tàu hỏa đang chuyển động với vận tốc $60 \ km/h$ thì bị hãm phanh chuyển động chậm dần đều. Sau khi đi thêm được $450\ m$ thì vận tốc của tàu chỉ còn $15 \ km/h$. Quãng đường tàu còn đi thêm được đến khi dừng hẳn là
\choice
{$60\ m$}
{$45\ m$}
{$15\ m$}
{\True $30\ m$}
\loigiai{}
\end{ex}
\begin{ex}%Tailieuchuan
Một ô tô chuyển động chậm dần đều. Sau $10\ s$, vận tốc của ô tô giảm từ $6 \ m/s$ về $4 \ m/s$. Quãng đường ô tô đi được trong khoảng thời gian $10\ s$ đó là
\choice
{$70\ m$}
{\True $50\ m$}
{$40\ m$}
{$100\ m$}
\loigiai{}
\end{ex}
\begin{ex}%Tailieuchuan
Một ô tô đang chạy với tốc độ $10 \ m/s$ trên đoạn đường thẳng thì người lái xe hãm phanh và ô tô chuyển động chậm dần đều. Cho tới khi dừng hẳn thì ô tô đã chạy thêm được $100\ m$. Gia tốc của xe có độ lớn bằng
\choice
{\True $0,5 \ m/s^2$}
{$0,2 \ m/s^2$}
{$0,7 \ m/s^2$}
{$0,9 \ m/s^2$}
\loigiai{

}
\end{ex}
\begin{ex}%Tailieuchuan
Một người đi xe đạp lên một cái dốc dài $50\ m$, chuyển động chậm dần đều với vận tốc lúc bắt đầu lên dốc là $18 \ km/h$, vận tốc ở đỉnh dốc là $3 \ m/s$. Gia tốc của xe có độ lớn là
\choice
{$16 \ m/s^2$}
{\True $0,16 \ m/s^2$}
{$1,6 \ m/s^2$}
{$0,26\ m/s^2$}
\loigiai{}
\end{ex}
\begin{ex}%Tailieuchuan
Xe chạy chậm dần đều lên một cái dốc dài $50\ m$, tốc độ ở chân dốc là $54 \ km/h$, ở đỉnh dốc là $36 \ km/h$. Sau khi lên được nửa dốc thì tốc độ của xe bằng
\choice
{$11,32 \ m/s$}
{$12,25 \ m/s$}
{\True $12,75 \ m/s$}
{$13,35 \ m/s$}
\loigiai{

}
\end{ex}
\begin{ex}%Tailieuchuan
Một ô tô đang chuyển động với vận tốc $10\ m/s$ thì bắt đầu tăng tốc, chuyển động nhanh dần đều. Sau $20\ s$ ô tô đạt được vận tốc $14 \ m/s$. Sau $50\ s$ kể từ lúc tăng tốc, gia tốc và vận tốc của ô tô lần lượt là
\choice
{$0,2 \ m/s^2$ và $18 \ m/s$}
{\True $0,2 \ m/s^2$ và $20 \ m/s$}
{$0,4 \ m/s^2$ và $38 \ m/s$}
{$0,1 \ m/s^2$ và $28 \ m/s$}
\loigiai{

}
\end{ex}









\subsubsection*{PHẦN 2. Thí sinh trả lời từ câu 1 đến câu 4. Trong mỗi ý a), b), c), d) ở mỗi câu, thí sinh chọn đúng hoặc sai.}
\setcounter{ex}{0}
\begin{ex}%book%8.10
Một ô tô tải đang chạy trên đường thẳng với vận tốc $18\ km/h$ thì tăng dần đều vận tốc. Sau $20\ s$, ô tô đạt được vận tốc $36\ km/h$.
\choiceTF[t]
{\True Chuyển động của ô tô là nhanh dần đều}
{Gia tốc của ô tô là $0,75\ m/s^2$}
{\True Tốc độ của ô tô đạt được sau $40\ s$ là $15\ m/s$}
{\True Sau $60\ s$ kể từ khi tăng tốc, ô tô đạt vận tốc $72\ km/h$}
\haidong{6}
\loigiai{
\begin{itemchoice}
\itemch 
\itemch $a=\dfrac{\Delta v}{\Delta t}=\dfrac{10-5}{20}=0,25\ m/s^2$.
\itemch $v_1=v_0+at_1=5+0,25.40=15\ m/s$.
\itemch $v_2=v_0+at_2\Rightarrow t_2=\dfrac{v_2-v_0}{a}=\dfrac{20-5}{0,25}=60\ s$.
\end{itemchoice}	
}
\end{ex}
\begin{ex}%book
\immini[thm]{Đồ thị ở hình bên mô tả sự thay đổi vận tốc theo thời gian trong chuyển động của một ô tô thể thao đang chạy thử về phía Bắc.}{\begin{tikzpicture}[scale=1.2,thick,>=stealth',scale=1.3,draw=Mapcolor]
\tkzDefPoints{0/0/o, 0/-1.2/i, 4/0/t, 0/1.5/d, 0.5/1/a, 1.5/1/b, 3.5/-1/c, 5/-0.5/e}
\draw  [help  lines, xstep=0.5cm,ystep=0.5cm]  (0,-1)  grid  (3.5,1);
\tkzDrawSegments[->](o,t i,d)
\tkzDrawSegments[very thick,Mapcolor](o,a a,b b,c)
\tkzLabelPoint[above ](t){$t\ (s)$}
\tkzLabelPoint[left](d){$v\ (m/s)$}
\tkzLabelPoint[left](o){$O$}
\foreach  \x  [evaluate=\x as \z using int(8*\x)] in  {0.5,1,...,3.5}
\node[below] at(\x,0){$\z$};
\foreach  \x  [evaluate=\x as \z using int(20*\x)] in  {-1,-0.5,0.5,1}
\node[left] at(0,\x){$\z$};
\end{tikzpicture}} 
\choiceTF[t]			
{Gia tốc của ô tô trong $4\ s$ đầu là $10\ m/s^2$}
{\True Gia tốc của ô tô trong từ giây thứ $4$ đến giây thứ $12$ có giá trị đại số là $0\ m/s^2$}
{\True Gia tốc của ô tô trong từ giây thứ $12$ đến giây thứ $20$ có giá trị đại số là $-2,5\ m/s^2$}
{\True Gia tốc của ô tô trong từ giây thứ $20$ đến giây thứ $28$ có giá trị đại số là $-2,5\ m/s^2$}
\loigiai{
\begin{itemchoice}
\itemch Trong 4 s đầu: $\Delta v=20\ m/s\ \Delta t=4\ s\Rightarrow a=\dfrac{\Delta v}{\Delta t}=\dfrac{20}{4}=5\ m/s^2$.
\itemch Từ giây thứ 4 đến giây thứ 12: $\Delta v=20-20=0\ m/s\ \Delta t=12-4=8\ s\Rightarrow a=\dfrac{\Delta v}{\Delta t}=0\ m/s^2$.
\itemch  Từ giây thứ 12 đến giây thứ 20: $\Delta v=0-20=-20\ m/s\ \Delta t=20-12=8\ s \Rightarrow a=\dfrac{\Delta v}{\Delta t}=\dfrac{-20}{8}=-2,5\ m/s^2$.
\itemch Từ giây thứ 20 đến giây thứ 28: $\Delta v=-20-0=-20\ m/s\ \Delta t=28-20=8\ s\Rightarrow a=\dfrac{\Delta v}{\Delta t}=\dfrac{-20}{8}=-2,5\ m/s^2$
\end{itemchoice}
}
\haidong{10} \end{ex}
\begin{ex}%book
\immini[thm]{Quan sát đồ thị $(v-t)$ trong hình bên của một vật đang chuyển động thẳng.}{\begin{tikzpicture}[scale=1,thick,>=stealth',draw=Mapcolor]
\tkzDefPoints{0/0/o, 2.75/0/x, 0/2.5/y}
\draw  [dashed,help  lines, xstep=0.5cm,ystep=0.5cm]  (0,0)  grid  (2,2);
\tkzDrawSegments[->](o,x o,y)
\tkzLabelPoint[below left](o){$O$}
\tkzLabelPoint[below](x){$t\ (s)$}
\tkzLabelPoint[above](y){$v\ (m/s)$}
\foreach  \x  in  {1,2,...,4}
{\pgfmathtruncatemacro{\label}{\x};
\node[below]
(\label) at (.5*\x,0) {\label};}
\foreach  \y  in  {2,4}
{\pgfmathtruncatemacro{\label}{\y};
\node[left]
(\label) at (0,.5*\y) {\label};}
\draw[very thick, Mapcolor] (0,0) -- (0.5,2) -- (1,2) -- (1.5,1) -- (2,1);
\end{tikzpicture}}
\choiceTF[t]
{Trong $2\ s$ đầu vật chuyển động thẳng nhanh dần đều}
{\True Vận tốc lớn nhất mà vật đạt được là $4\ m/s$}
{Gia tốc của vật trong khoảng thời gian từ $0\ s$ đến $1\ s$ là $16\ m/s^2$}
{\True Quãng đường vật đi được trong khoảng thời gian từ $1\ s$ đến $2\ s$ lớn hơn quãng đường đi được trong khoảng thời gian từ $0$ đến $1\ s$}
\loigiai{\begin{itemchoice}
\itemch Trong khoảng thời gian từ $0$ đến $1\ s$ vật chuyển động thẳng nhanh dần đều. Trong khoảng thời gian từ $1 s$ đến $2 s$ vật chuyển động thẳng đều
\itemch
\itemch Gia tốc của vật trong khoảng thời gian từ 0 s đến 1 s là $a=\dfrac{v-v_{0}}{t}=\dfrac{4-0}{1}=4\ m/s^{2}$
\itemch $S_{0 \rightarrow 1\ s}=0,5 . S_{1 \rightarrow 2\ s}$
\end{itemchoice}
}

\haidong{10} \end{ex}
\begin{ex}%book%SBT9.11
Hai vật $A$ và $B$ chuyển động cùng chiều trên đường thẳng có đồ thị vận tốc - thời gian vẽ ở hình bên. Biết ban đầu hai vật cách nhau $78\ m$.
\immini{\choiceTF[t]			
{Phương trình vận tốc của vật $A$ là $v=40-t\ (m\ s)$}
{\True Phương trình vận tốc của vật $B$ là $v=t\ (m\ s)$}
{\True Hai vật có cùng vận tốc ở thời điểm xấp xỉ $13,3 \ s$}
{Hai vật gặp nhau cách vị trí ban đầu của $A$ là $81,5\ m$}
}{\begin{tikzpicture}[scale=1,thick,>=stealth',x=2cm,draw=Mapcolor]
\tkzDefPoints{0/0/o, 2.5/0/t, 0/5/d}
\draw  [help  lines, xstep=2cm,ystep=1cm]  (0,0)  grid  (2,4);
\tkzDrawSegments[->](o,t o,d)
\tkzLabelPoint[below](t){$t\ (s)$}
\tkzLabelPoint[left](d){$v\ (m/s)$}
\tkzLabelPoint[below left](o){$O$}
\foreach  \x  [evaluate=\x as \z using int(10*\x)] in  {1,2}
\node[below] at(\x,0){$\z$};
\foreach  \x  [evaluate=\x as \z using int(10*\x)] in  {1,2,...,4}
\node[left] at(0,\x){$\z$};
\draw[very thick, Mapcolor](0,0)--(2,2)node[above right,black]{$B$};
\draw[very thick, Mapcolor](0,4)--(2,0)node[above right,black]{$A$};
\end{tikzpicture}}
\haidong{10}
\loigiai{
\begin{itemchoice}
\itemch 
\itemch 
\itemch $v_1=v_2$ tại thời điểm $t\approx 13,3\ s$.
\itemch Chuyển động A: $v_{01}=40\ m/s\ a_1=\dfrac{0-40}{20}=-2\ m/s^2\ d_1=40t-t^2\quad (1)$.\\
Chuyển động B: $v_{02}=0\ m/s\ a_2=\dfrac{10-0}{10}=1\ m/s^2\ d_1=78+0,5t^2$.
\itemch Hai vật gặp nhau khi $d_1=d_2\Rightarrow 40t-t^2=78+0,5t^2$

Phương trình cho hai nghiệm $t=2,12\ s$ và $t'=24,5\ s$.

Loại $t'$ vì $t'>20\ s$.

Thay t vào $(1)$ ta được $d_1=81,5\ m$.

Hai vật gặp nhau cách vị trí ban đầu của A là 81,5 m.
\end{itemchoice}
}
\haidong{10} \end{ex}
\subsubsection*{PHẦN 3. Thí sinh trả lời từ câu 1 đến câu 6.}
\setcounter{ex}{0}
\noindent\textit{Sử dụng thông tin sau cho Câu 1 và Câu 2:} Một máy bay có vận tốc khi tiếp đất là $100\ m/s$. Để giảm vận tốc sau khi tiếp đất, máy bay chỉ có thể có gia tốc đạt độ lớn cực đại là $4\ m/s^2$.
\begin{ex}%book%SBT-KNTT9.7
Thời gian ngắn nhất để máy bay dừng hẳn kể từ khi tiếp đất là bao nhiêu giây?
\shortans[oly]{25}
\loigiai{
\begin{itemize}
\item $v_0=100\ m/s\ a=-4\ m/s^2\ v=0$.
\item $v=v_0+at\Rightarrow t=\dfrac{v-v_0}{a}=25\ s$.
\end{itemize}
}
\haidong{7}
\end{ex}
\begin{ex}		
Máy bay này có thể hạ cánh an toàn ở sân bay có đường bay dài $1\ km$ hay không? Nếu có điền số $1$; nếu sai điền số $0$.
\shortans[oly]{0}
\haidong{7}
\loigiai{
\begin{itemize}
\item $d=v_0t+\dfrac{at^2}{2}=1250\ m=1,25\ km$.
\item			Máy bay không thể hạ cánh an toàn trên sân bay có đường bay dài 1 km.
\end{itemize}
}
\haidong{7} \end{ex}











\textit{Sử dụng thông tin sau cho Câu 3, Câu 4 và Câu 5:} Quan sát đồ thị $(v-t)$ mô tả chuyển động thẳng của tàu hoả trong hình.
\begin{center} 
\begin{tikzpicture}[line cap=round, line join=round, scale=1, font=\footnotesize, >=stealth,draw=Mapcolor]
\path
(0,0)coordinate(O)
(6.5,0)coordinate(x)
(0,2.5)coordinate(y)
;
\draw  [help  lines, xstep=0.5cm,ystep=0.5cm]  (0,0)  grid  (6.5,2);
\draw [->](O)--(x);
\draw [->](O)--(y);

\node at (0,3) {$v\ (m/s)$};
\node at (7,0) {$t\ (s)$};
\node[below left] at (0,0) {$O$};

\foreach  \x  in  {1,2,...,12}
{\pgfmathtruncatemacro{\label}{10*\x};
\node[below]
(\label) at (.5*\x,0) {\label};
\draw[fill=white](.5*\x,0)circle(0.5pt);}


\foreach  \y  in  {1,2,3,4}
{\pgfmathtruncatemacro{\label}{10*\y};
\node[left]
(\label) at (0,.5*\y) {\label};
\draw[fill=white](0,.5*\y)circle(0.5pt);}

\draw[very thick,Mapcolor] (0,0.5) .. controls (3.2,2.2) and (2.2,2.2) .. (4.5,0);
\draw[very thick,Mapcolor] (4.5,0) -- (5.5,0);
\draw[very thick,Mapcolor] (5.5,0) .. controls (5.6,0.6) and (6,1) .. (6.48,1);
\end{tikzpicture}
\end{center}
\begin{ex}%Tailieuchuan
Tại thời điểm bao nhiêu giây thì vận tốc tàu hoả có giá trị lớn nhất?
\shortans[oly]{50}
\loigiai{
}
\end{ex}
\begin{ex}%Tailieuchuan
Tàu hỏa chuyển động với vận tốc ban đầu bằng bao nhiêu $m/s$?
\shortans[oly]{10}
\loigiai{
}
\end{ex}
\begin{ex}%Tailieuchuan
Tàu hỏa chuyển động với gia tốc bao nhiêu $m/s^2$ trong $50\ s$ đầu tiên?
\shortans[oly]{0,5}
\end{ex}
\begin{ex}%Tailieuchuan
Một người đi xe máy đang chuyển động với vận tốc $10\ m/s$. Để không va vào con chó, người ấy phanh xe. Biết độ dài vết phanh xe là $5\ m$. Người đi xe máy chuyển động với gia tốc có độ lớn bao nhiêu $m/s^2$?
\shortans[oly]{10}
\end{ex}
\newpage 

\TieudeT{PHONG TỎA VẬT LÍ}
\subsubsection*{PHẦN I. Thí sinh trả lời câu hỏi từ câu 1 đến câu 18. Mỗi câu hỏi thí sinh chỉ chọn một phương án}
\setcounter{ex}{0}
\begin{ex}%book%[1P1-5.2K][Loai1]
Một đoàn tàu đang chạy với tốc độ $54\ km/h$ thì hãm phanh. Sau đó chạy thêm được $25\ m$ thì dừng hẳn. Độ lớn gia tốc của đoàn tàu là
\choice
{$3\ m/s^2$}
{\True $4,5\ m/s^2$}
{$5\ m/s^2$}
{$6\ m/s^2$}
\haidong{4}
\loigiai{
\begin{itemize}
\item $54\ km/h = 15\ m/s$.
\item Áp dụng công thức độc lập thời gian ta có:
\begin{align*}
v^2-v_0^2=2ad \Rightarrow a=\dfrac{v^2-v_0^2}{2d}=\dfrac{0-15^2}{2,25}=-4,5\ m/s^2.
\end{align*}	
\end{itemize}
}
\haidong{7} 
\end{ex}
\begin{ex}%book%[1P1-5.2K][Loai1]
Một xe đang chuyển động với tốc độ $36\ km/h$ thì hãm phanh sau đó chuyển động chậm dần đều, sau $20$ giây tốc độ là $18\ km/h$. Sau bao lâu kể từ lúc hãm phanh thì xe dừng lại?
\choice
{$30\ s$}
{\True $40\ s$}
{$42\ s$}
{$50\ s$}	
\haidong{5}
\loigiai{
\begin{itemize}
\item Đổi $36\ km/h = 10\ m/s\ 18\ km/h = 5\ m/s.$
\item Ta có: $v =v_0+at \Rightarrow 5=10+20a \Rightarrow a=-0,25\ m/s^2$.  
\item Vận tốc của xe lúc dừng lại là: $v=0\ m/s$. 
\item Thời gian từ lúc hãm phanh đến lúc xe dừng lại là: $t=\dfrac{v-v_0}{a}=\dfrac{0-10}{-0,25}=40\ s$. 
\end{itemize}	
}
\haidong{7} \end{ex}
\begin{ex}%book%[1P1K5-2][Loai1]
Một đoàn tàu đang chạy với tốc độ $72\ km/h$ thì hãm phanh, chạy thẳng chậm dần đều, sau $10\ s$ tốc độ giảm xuống còn $54\ km/h$. Sau bao lâu kể từ lúc hãm thì tàu dừng lại?
\choice
{$t = 30\ s$}
{$t = 60\ s$}
{\True $t = 40\ s$}
{$t = 50\ s$}	
\haidong{4}
\loigiai{
\begin{itemize}
\item $72\ km/h = 20\ m/s\ 54\ km/h = 15\ m/s$.
\item Ta có: $v_t=v_0+at$. 
\item Gia tốc của tàu là: $a=\dfrac{v_t-v_0}{t}=-0,5\ m/s^2$. 
\item Khi tàu dừng lại thì $v_t = 0\ m/s$.   
\item Thời gian từ lúc hãm đễn lúc tàu dừng lại là:  $t=\dfrac{v_t-v_0}{a}=\dfrac{0-20}{-0,5}=40\ s$.
\end{itemize}	
}
\haidong{7} \end{ex}
\begin{ex}%book
Một ô tô chuyển động chậm dần đều. Sau $10\ s$, vận tốc của ô tô giảm từ $6\ m/s$ về $4\ m/s$. Quãng đường ô tô đi được trong khoảng thời gian $10\ s$ đó là
\choice
{$40\ m$}
{$100\ m$}
{$70\ m$}
{\True $50\ m$}
\haidong{5}
\loigiai{
}
\haidong{7} \end{ex}
\begin{ex}%book%[1P1K5-3][Loai1]
Vận tốc của một chất điểm chuyển động dọc theo trục $Ox$ cho bởi hệ thức: $v_t=10-2t\ (m/s)$. Vận tốc trung bình của chất điểm trong khoảng thời gian từ $t_1 = 2\ s$ đến $t_2 = 4\ s$ là
\choice
{$7\ m/s$}
{$6\ m/s$}
{$5\ m/s$}
{\True $4\ m/s$}
\haidong{7}
\loigiai{
\begin{itemize}
\item Ta có:
\begin{align*}
v_{tb}=\dfrac{d_2-d_1}{t_2-t_1}& =\dfrac{\left(v_0t_2+0,5at_2^2\right)-\left(v_0t_1+0,5at_1^2\right)}{t_2-t_1}\\
\Leftrightarrow v_{tb}& =\dfrac{v_0\left(t_2-t_1\right)+0,5a\left(t_2^2-t_1^2\right)}{t_2-t_1}\\
\Rightarrow v_{tb}& =v_0+\dfrac{1}{2}a\left(t_2+t_1\right)=\dfrac{v_0+at_1}{2}+\dfrac{v_0+at_2}{2}=\dfrac{v_1+v_2}{2}.
\end{align*}
\item Mà: $\begin{cases}
v_1=10-2t_1=10-2.2=6\ m/s.\\
v_2=10-2t_2=10-2.4=2\ m/s.
\end{cases}\Rightarrow v_{tb}=4\ m/s$.
\end{itemize}
}
\haidong{7} \end{ex}
\begin{ex}%book%[1P1G5-3][Loai1] 
Một ô tô đang chuyển động thẳng đều với tốc độ $72\ km/h$ thì giảm đều tốc độ, khi đi được quãng đường $75\ m$ thì tốc độ chỉ còn lại một nửa ban đầu. Quãng đường đi được trong $2$ giây cuối trước khi dừng hẳn là
\choice
{\True $4\ m$}
{$5\ m$}
{$6\ m$}
{$7\ m$}
\haidong{8}
\loigiai{
\begin{itemize}
\item Ta có: $\begin{cases}
v_t^2-v_0^2=2as.\\
v_t=\dfrac{v_0}{2}\ s=50\ m\ v_0=20\ m/s.
\end{cases}\Rightarrow a=\dfrac{v_t^2-v_0^2}{2s}=-2\ m/s^2$.
\item Khi dừng hẳn thì: $v_t=0\Leftrightarrow 0=20-2t\Rightarrow t=10\ s$
\item Vậy: $s=v_0t+\dfrac{1}{2}at^2\Rightarrow \Delta s=s_{10}-s_8=20\left(10-8\right)-\dfrac{1}{2}.2.\left(10^2-8^2\right)=4\ m$.
\end{itemize}
}
\haidong{7} \end{ex}
\begin{ex}%book
\immini[thm]{Đồ thị vận tốc - thời gian $(v-t)$ của chuyển động thẳng được cho như hình vẽ bên. Chuyển động thẳng nhanh dần đều ứng với đoạn đồ thị nào?
\choice
{$MN$}
{NO}
{\True PQ}
{OP}}{\begin{tikzpicture}[very thick,draw=Mapcolor]
\tkzDefPoints{0/0/x, 0/3.5/v, 5.5/0/t, 0/2/m, 1/2/n, 2/0.5/o, 3/0.5/p, 5/1.5/q}
\tkzDrawSegments[->](x,t x,v)
\tkzDrawSegments[Mapcolor,very thick](m,n n,o o,p p,q)
\draw  [help  lines, dashed, xstep=0.5cm,ystep=0.5cm]  (0,0)  grid  (5.4,3.2);
\tkzLabelPoint[above left](m){M}
\tkzLabelPoint[above right](n){N}
\tkzLabelPoint[below](o){O}
\tkzLabelPoint[below](p){P}
\tkzLabelPoint[above right](q){Q}
\tkzLabelPoint[above](v){$v\ (m/s)$}
\tkzLabelPoint[below](t){$t~(s)$}
\tkzLabelPoint[below](x){$0$}
\end{tikzpicture}}

\loigiai{
}
\haidong{10} \end{ex}
\begin{ex}%book
\immini[thm]{Chất điểm chuyển động có đồ thị vận tốc theo thời gian như hình dưới đây. Tính độ dịch chuyển của người này từ khi bắt đầu chuyển động đến thời điểm $5$ giây.
\choice
{$12\ m$}
{\True $14\ m$}
{$16\ m$}
{$10\ m$}}{\begin{hinhanh}{6}\begin{tikzpicture}[very thick,draw=Mapcolor]
\tkzDefPoints{0/0/o, 0/3.5/v, 6/0/t, 1.5/2/a, 3.5/2/d, 5/3/f}
\tkzDrawSegments[->](o,v o,t)
\draw  [help  lines, dashed, xstep=0.5cm,ystep=0.5cm]  (0,0)  grid  (5.5,3);
\tkzDrawSegments[very thick, Mapcolor](o,a a,d d,f)
\tkzLabelPoint[above](v){$v~(m/s)$}
\tkzLabelPoint[below](t){$t~(s)$}
\tkzLabelPoint[below](o){$O$}
%\tkzLabelPoint[above left](a){$A$}
%\tkzLabelPoint[above left](d){$D$}
%\tkzLabelPoint[above left](f){$F$}
\node[left] at (0,1) {2};
\node [left] at (0,2) {4};
\node[left]  at (0,3) {6};
%\node [above left] at (2.5,2) {B};
%\node[above left] at (2.5,0) {C};
%\node [above left]at (3.5,0) {E};
%\node [above left]at (5.5,0) {G};
\node[below] at (1,0) {2};
\node[below] at (2,0) {4};
\node[below] at (3,0) {6};
\node[below] at (4,0) {8};
\node[below] at (5,0) {10};
\end{tikzpicture}\end{hinhanh}}
\haidong{3}

\loigiai{
}
\haidong{10} \end{ex}
\begin{ex}%book
\immini[thm]{Hình bên là đồ thị vận tốc của vật chuyển động thẳng theo ba giai đoạn liên tiếp. Tính chất chuyển động của vật trên đoạn $OA$ là
\choice
{chuyển động nhanh dần đều với gia tốc $a=12\ cm/s^2$}
{\True chuyển động nhanh dần đều với gia tốc $a=12\ m/s^2$}
{đứng yên}
{chuyển động chậm dần đều với gia tốc $a=-12\ m/s^2$}}{\begin{hinhanh}{6}\begin{tikzpicture}[very thick, scale=0.8,draw=Mapcolor]
\tkzDefPoints{0/0/0, 0/3/v, 5/0/t, 0/1/6, 0/2/12, 1/0/1, 2/0/2, 3/0/3, 4/0/4, 1/2/a, 3/2/b, 4/1/c}
\tkzDrawSegments[->, thick](0,v 0,t)
\tkzDrawSegments[very thick,Mapcolor](0,a a,b b,c)
\tkzDrawSegments[dashed, thin](12,a a,1 b,3 6,c c,4)
\tkzDrawPoints[size=1.5,fill=black](6,12,1,2,3,4)
\tkzLabelPoint[above](a){$A$}
\tkzLabelPoint[above](b){$B$}
\tkzLabelPoint[above right](c){$C$}
\tkzLabelPoint[left](6){$6$}
\tkzLabelPoint[below](1){$1$}
\tkzLabelPoint[below](2){$2$}
\tkzLabelPoint[below](3){$3$}
\tkzLabelPoint[below](4){$4$}
\tkzLabelPoint[left](12){$12$}
\tkzLabelPoint[below](0){0}
\tkzLabelPoint[above](v){$v~(m/s)$}
\tkzLabelPoint[below](t){$t~(s)$}
\end{tikzpicture}\end{hinhanh}}
\haidong{4}

\loigiai{
}
\haidong{10} \end{ex}
\begin{ex}%Tailieuchuan
Lúc $1\ h$, một xe qua $A$ với tốc độ $10 \ m/s$, chuyển động nhanh dần đều với gia tốc $1\ m/s^2$ đuổi theo một xe đạp đang chuyển động nhanh dần đều qua $B$ với tốc độ đầu là $2 \ m/s$ và với gia tốc là $0,5 \ m/s^2$. Sau $20\ s$ thì xe đuổi kịp xe đạp. Khoảng cách $AB$ là
\choice
{$300\ m$}
{$250\ m$}
{$200\ m$}
{\True $260\ m$}
\loigiai{}
\end{ex}
\begin{ex}%Tailieuchuan
Một ô tô đang chuyển động với vận tốc ban đầu là $10 \ m/s$ trên đoạn đường thẳng, thì người lái xe hãm phanh, xe chuyển động chậm dần với gia tốc $2 \ m/s^2$. Quãng đường mà ô tô đi được sau thời gian $3$ giây là
\choice
{$19 \ m$}
{$20 \ m$}
{$18 \ m$}
{\True $21 \ m$}
\loigiai{}
\end{ex}
\begin{ex}%Tailieuchuan
Một xe lửa bắt đầu rời khỏi ga và chuyển động thẳng nhanh dần đều với gia tốc $0,1 \ m/s^2$. Khoảng thời gian để xe đạt được vận tốc $36 \ km/h$ là
\choice
{$360 \ s$}
{$200 \ s$}
{$300 \ s$}
{\True $100 \ s$}
\loigiai{}
\end{ex}
\begin{ex}%Tailieuchuan
Một ôtô chuyển động thẳng nhanh dần đều. Sau $10\ s$, vận tốc của ô tô tăng từ $4\ m/s$ đến $6 \ m/s$. Quãng đường mà ôtô đã đi được trong khoảng thời gian này là
\choice
{$100 \ m$}
{\True $50 \ m$}
{$25\ m$}
{$500\ m$}
\loigiai{}
\end{ex}
\begin{ex}%Tailieuchuan
Một ô tô đang chuyển động với vận tốc $54 \ km/h$ thì người lái xe hãm phanh. Ô tô chuyển động thẳng chậm dần đều và sau $6$ giây thì dừng lại. Quãng đường $s$ mà ô tô chạy thêm được kể từ lúc hãm phanh là
\choice
{\True $45 \ m$}
{$82,6 \ m$}
{$252 \ m$}
{$135\ m$}
\loigiai{}
\end{ex}










\textit{Sử dụng thông tin sau cho Câu 17 và Câu 18:} Hai người đi xe đạp khởi hành cùng lúc và đi ngược chiều nhau. Người thứ nhất có vận tốc đầu là $18 \ km/h$ và chuyển động chậm dần đều với gia tốc có độ lớn $20\ cm/s^2$. Người thứ hai có vận tốc đầu là $5,4 \ km/h$ và chuyển động nhanh đều với gia tốc có độ lớn $0,2 \ m/s^2$. Khoảng cách giữa hai người là $130\ m$.
\begin{ex}%Tailieuchuan
Sau bao lâu thì hai người gặp nhau? 
\choice
{\True $20 \ s$}
{$18,2 \ s$}
{$27\ s$}
{$17,5 \ s$}
\loigiai{}
\end{ex}
\begin{ex}%Tailieuchuan
Vị trí gặp nhau cách 
\choice
{\True $A$ $60\ m$}
{$A$ $56,9\ m$}
{$B$ $60\ km$}
{$B$ $56,9\ m$}
\loigiai{}
\end{ex}
\begin{ex}%Tailieuchuan
Khi ô tô đang chạy với vận tốc $10 \ m/s$ trên đoạn đường thẳng thì người lái xe hãm phanh và ô tô chuyển động chậm dần đều. Cho tới khi dứng hẳn lại thì ô tô đã chạy thêm được $100\ m$. Gia tốc của ô tô có độ lớn là
\choice
{\True $0,5 \ m/s^2$}
{$a=0,2 \ m/s^2$}
{$a=0,4 \ m/s^2$}
{$a=0,1\ m/s^2$}
\loigiai{}
\end{ex}
\begin{ex}%Tailieuchuan
Thời gian cần thiết để tăng vận tốc từ $10 \ m/s$ lên $40 \ m/s$ của một chuyển động có gia tốc $2 \ m/s^2$ là
\choice
{$10\ s$}
{\True $15\ s$}
{$25\ s$}
{$20\ s$}
\loigiai{}
\end{ex}

\subsubsection*{PHẦN 2. Thí sinh trả lời từ câu 1 đến câu 4. Trong mỗi ý a), b), c), d) ở mỗi câu, thí sinh chọn đúng hoặc sai.}
\setcounter{ex}{0}
\begin{ex}%Tailieuchuan
\immini[thm]{Chuyển động của một xe được mô tả bởi đồ thị sau.
\choiceTF[t]
{\True Tốc độ của xe có giá trị lớn nhất tại thời điểm $t=0$} 
{Xe chuyển động nhanh dần đêu trong $20\ s$ đầu} 
{\True Gia tốc của xe trong khoảng $80\ s$ cuối không thay đổi} 
{\True Tổng quãng đường xe đi được là $1\ km$}
}{\begin{tikzpicture}[line cap=round, line join=round, scale=1, font=\footnotesize, >=stealth,draw=Mapcolor]
\path
(0,0)coordinate(O)
(6.5,0)coordinate(x)
(0,5.5)coordinate(y)
;
\draw [->](O)--(x);
\draw [->](O)--(y);

\node at (0,6) {$v\ (m/s)$};
\node at (7,0) {$t\ (s)$};
\node[below left] at (0,0) {$O$};

\draw[dashed] (0,2) -- (1,2) -- (1,0);
\draw[dashed] (2.5,0) -- (2.5,2);
\draw[very thick,Mapcolor] (0,4) -- (1,2) -- (2.5,2) -- (6,0);
\node[left] at (0,2) {$10$};
\node[left] at (0,4) {$20$};
\node[below] at (1,0) {$20$};
\node[below] at (2.5,0) {$50$};
\node[below] at (6,0) {$130$};
\end{tikzpicture}} 
\loigiai{ 
\begin{itemchoice}
\itemch
\itemch
\itemch
\itemch
\end{itemchoice}
}
\end{ex}
\begin{ex}%book%SBT-KNTT9.9
Một người đi xe máy trên một đoạn đường thẳng muốn đạt được vận tốc $36\ km/h$ sau khi đi được $100\ m$ bằng một trong hai cách sau:
\begin{itemize}
\item \textit{Cách 1:} Chạy thẳng nhanh dần đều trong suốt quãng đường.
\item \textit{Cách 2:} Chỉ cho xe chạy nhanh dần đều trên $\dfrac{1}{5}$ quãng đường, sau đó cho xe chuyển động thẳng đều trên quãng đường còn lại.
\end{itemize}
\choiceTF[t]
{\True Cách $1$ đi mất $20\ s$}
{Thời gian chuyển động nhanh dần đều trong cách thứ $2$ là $8\ s$}
{Thời gian chuyển động thẳng đều trong cách thứ $2$ là $4\ s$}
{\True Cách $2$ đi mất $12\ s$}
\haidong{9}
\loigiai{
\begin{itemchoice}
\itemch \textbf{Cách 1:} Dùng công thức  $v^2-v_0^2=2ad$.\\
Vì $v_0=0\ v=10\ m/s\ d=100\ m$ nên $a=0,5\ m/s^2$.\\
Vì $v=v_0+at\Rightarrow t=20\ s$.
\itemch \textbf{Cách 2:} Thời gian chuyển động nhanh dần đều.\\
Vì $v_0=0\ v=10\ m/s\ d_1=\dfrac{100}{5}=20\ m$ nên $a_1=2,5\ m/s^2$ và $t_1=4\ s$.
\itemch Thời gian chuyển động đều: $t_2=\dfrac{100-d_1}{v}=\dfrac{80}{10}=8\ s$.\\
\itemch Thời gian chuyển động trong cách 2: $t'=t_1+t_2=12\ s$.
\end{itemchoice}
}
\haidong{7} \end{ex}
\begin{ex}%book
\immini[thm]{Hình bên là đồ thị vận tốc - thời gian trong chuyển động thẳng của một bạn đang đi trong siêu thị theo hướng Bắc.}{\begin{tikzpicture}[scale=1,thick,>=stealth',draw=Mapcolor]
\tkzDefPoints{0/0/o, 0/-1/i, 5.5/0/t, 0/2.5/d, 0/1.5/a, 2/1.5/b, 3/0/c, 3.5/0/e, 4/-0.5/f, 4.5/-0.5/g, 5/0/h}
\draw  [help  lines, xstep=0.5cm,ystep=0.5cm]  (0,-1)  grid  (5,2);
\tkzDrawSegments[->](o,t i,d)
\tkzDrawSegments[line width=1.2pt,Mapcolor](a,b b,c c,e e,f f,g g,h)
\tkzLabelPoint[left](o){$O$}
\node[rotate=90,yshift=30pt] at (0,0.5){Vận tốc (m/s)};
\node[above] at (4.5,0){Thời gian (s)};
\foreach  \x  [evaluate=\x as \z using int(2*\x)] in  {0.5,1,...,5}
\node[below] at(\x,0){$\z$};
\foreach  \x  [evaluate=\x as \z using int(1*\x)] in  {-1,1}
\node[left] at(0,\x){$\z$};
\node[left] at(0,-0.5){$-0,5$};
\node[left] at(0,0.5){$0,5$};
\node[left] at(0,1.5){$1,5$};
\end{tikzpicture}}
\choiceTF[t]
{\True Trong $4\ s$ đầu tiên: bạn đó chuyển động đều với tốc độ $1,5\ m/s$}
{\True Từ giây $6\ s$ đến $7\ s$: bạn đó đứng yên}
{\True Từ $7\ s$ đến $8\ s$: bạn đó quay lại đi theo hướng Nam}
{Từ $8\ s$ đến $9\ s$: bạn đó chuyển động thẳng đều theo hướng Bắc với tốc độ $0,5\ m/s$}
\loigiai{
\begin{itemchoice}
\itemch Trong 4 s đầu tiên: bạn đó đi đều với vận tốc 1,5 m/s.
\itemch Từ giây 6 đến giây 7: bạn đó nghỉ
\itemch Từ giây 7 đến giây 8: bạn đó bắt đầu đi theo chiều âm
\itemch Từ giây 8 – 9: bạn đó đi đều với vận tốc $-0,5$ m/s.
\end{itemchoice}
}
\haidong{10} \end{ex}
\begin{ex}%book
\immini[thm]{Một vật chuyển động có đồ thị biểu diễn vận tốc theo thời gian như hình bên.}{\begin{tikzpicture}[scale=1.2,thick,>=stealth',draw=Mapcolor]
\tkzDefPoints{0/0/o, 3.5/0/x, 0/2.5/y}
\tkzDrawSegments[->](o,x o,y)
\tkzLabelPoint[below left](o){$O$}
\tkzLabelPoint[above](x){$t\ (s)$}
\tkzLabelPoint[above](y){$v\ (m/s)$}
\draw[dashed,thin] (0,2) -- (1,2) -- (1,0);
\draw[very thick, Mapcolor] (0,0) -- (1,2) -- (3,0);
\node[left] at (0,2) {$30$};
\node[below] at (1,0) {$15$};
\node[below] at (3,0) {$40$};
\end{tikzpicture}}
\choiceTF[t]
{Trong thời gian $40\ s$, vật chuyên động chậm dần đều rồi nhanh dần đều}
{\True Gia tốc của vât trong $15\ s$ đầu là $2\ m/s^{2}$}
{Độ lớn gia tốc của vật từ $t=15\ s$ đến $t=40\ s$ là $0,75\ m/s^{2}$}
{Quãng đường vật đi được trong $15\ s$ đầu lớn hơn quãng đường vật đi được trong $25\ s$ sau}
\loigiai{\begin{itemchoice}
\itemch. Vật chuyển động nhanh dần đều rồi chậm dần đều. Ta thấy đường biểu diễn vận tốc theo thời gian là hai đoạn thẳng xiên góc nên đây là chuyển động biến đổi đều. Trong thời gian từ $t=0\ s$ đến $t=15\ s$ vận tốc luôn dương, vận tốc tăng suy ra vật chuyển động nhanh dần đều, trong khoảng thời gian từ $t=15\ s$ đến $t=40\ s$ vận tốc luôn dương, vận tốc giảm suy ra vật chuyển động chậm dần đều.
\itemch Gia tốc của vật trong 15 s đầu là $a=\dfrac{v-v_{0}}{t}=\dfrac{30-0}{15}=2\ \left(m / s^{2}\right)$
\itemch Độ lớn gia tốc của vật từ $t=15\ s$ đến $t=40\ s$ là $|a|=\left|\dfrac{v_{1}-v}{t^{\prime}}\right|=\left|\dfrac{0-30}{25}\right|=1,2\ \left(m / s^{2}\right)$
\itemch 
\begin{center}
\begin{tikzpicture}[scale=1,thick,>=stealth',draw=Mapcolor]
\tkzDefPoints{0/0/o, 3.75/0/x, 0/2.5/y}
\tkzDrawSegments[->](o,x o,y)
\tkzLabelPoint[below left](o){$O$}
\tkzLabelPoint[above](x){$t\ (s)$}
\tkzLabelPoint[above](y){$v\ (m/s)$}
\draw[dashed,thin] (0,2) -- (1,2) -- (1,0);
\draw[very thick, Mapcolor] (0,0) -- (1,2) -- (3,0);
\node[left] at (0,2) {$30$};
\node[below] at (1,0) {$15$};
\node[below] at (3,0) {$40$};
\node[above left] at (0,0) {$A$};
\node[above] at (1,2) {$B$};
\node[above right] at (1,0) {$H$};
\node[above] at (3,0) {$C$};
\end{tikzpicture}
\end{center} 
Quãng đường đi được trong 15 đầu bằng diện tích hình tam giác vuông ABH

Quãng đường đi được trong 25 s đầu bằng diện tích hình tam giác vuông BCH

Hai tam giác này có cùng đường cao nhưng cạnh đáy AH (tương ứng 15 s) nhỏ hơn cạnh đáy HC (tương ứng 25 s) nên tam giác ABH có diện tích nhỏ hơn tam giác BCH. Vậy quãng đường đi được trong 15 s đầu nhỏ hơn quãng đường đi được trong 25 s sau.
\end{itemchoice}
}
\haidong{10} \end{ex}


























\subsubsection*{PHẦN 3. Thí sinh trả lời từ câu 1 đến câu 6.}
\setcounter{ex}{0}
\textit{Sử dụng thông tin sau cho Câu 1 và Câu 2:} Một con báo đang chạy với vận tốc $30\ m/s$ thì chuyển động chậm dần đều khi tới gần một con suối. Trong $3$ giây, vận tốc của nó giảm còn $9\ m/s$.
\begin{ex}%Tailieuchuan
Con báo chạy với gia tốc có độ lớn là bao nhiêu $m/s^2$?
\shortans[oly]{7}
\loigiai{}
\end{ex}
\begin{ex}%Tailieuchuan
Quãng đường con báo chạy được sau $5\ s$ là bao nhiêu mét?
\shortans[oly]{62,5}
\loigiai{}
\end{ex}
\begin{ex}%Tailieuchuan
Xe thứ nhất đi từ $A$ đến $B$ trên một đường thẳng với vận tốc $10\ m/s$, chuyển động chậm dần đều với gia tốc có độ lớn $0,2\ m/s^2$. Cùng lúc đó tại một điểm $B$ cách $A$ $560\ m$, một xe thứ hai khởi hành đi ngược chiều với xe thứ nhất chuyển động nhanh dần đều với gia tốc $0,4\ m/s^2$. Hai xe gặp nhau sau bao nhiêu giây?
\shortans[oly]{40}
\loigiai{}
\end{ex}
\begin{ex}%Tailieuchuan
Một vật chuyển động thẳng nhanh dần đều có vận tốc đầu là $18\ km/h$. Trong giây thứ $5$, vật đi được quãng đường là $5,9\ m$. Quãng đường vật đi được sau khoảng thời gian $10\ s$ kể từ khi vật bắt đầu chuyển động là bao nhiêu cm?
\shortans[oly]{6000}
\loigiai{
}
\end{ex}
\begin{ex}%Tailieuchuan
Một ô tô chuyển động thẳng nhanh dần đều. Sau $10\ s$, vận tốc của ô tô tăng từ $4\ m/s$ đến $6\ m/s$. Quãng đường mà ô tô đi được trong khoảng thời gian trên là bao nhiêu km?
\shortans[oly]{0,05}
\loigiai{
}
\end{ex}
\begin{ex}%Tailieuchuan
Một xe máy đi qua $A$ với tốc độ $10\ m/s$, chuyển động nhanh dần đều với gia tốc $1\ m/s^2$ đuổi theo một xe đạp đang chuyển động chậm dần đều qua $B$ với tốc độ ban đầu là $3\ m/s$ và với độ lớn gia tốc là $0,5\ m/s^2$. Sau $10\ s$ thì xe đuổi kịp xe đạp. Khoảng cách $AB$ là bao nhiêu mét?
\shortans[oly]{145}
\loigiai{}
\end{ex}


































































